\newpage
\section{Tafeel}

% ─── Bab name & example ─────────────────────────────────────
\begin{center}
  \textbf{এই নতুন বাবের নাম}
  \enspace
  {\arabicfont \textarabic{بَابُ التَّفْعِيلِ}}
\end{center}

\vspace{0.3cm}

\begin{center}
  \textbf{শেখানো} \enspace---\enspace
  {\arabicfont \textarabic{التَّعْلِيمُ}}
\end{center}

\vspace{0.6cm}

% ─── Madi forms ─────────────────────────────────────────────
\begin{center}
  {\arabicfont \Large
    \textarabic{عَلَّمَ - عَلَّمَتْ - عَلَّمْتَ - عَلَّمْتِ - عَلَّمْتُ}
  }
\end{center}

\vspace{0.4cm}

% ─── Mudare forms ───────────────────────────────────────────
\begin{center}
  {\arabicfont \Large
    \textarabic{يُعَلِّمُ - تُعَلِّمُ - تُعَلِّمُ - تُعَلِّمِينَ - أُعَلِّمُ}
  }
\end{center}

\vspace{0.4cm}

% ─── Amr / Nahi forms ───────────────────────────────────────
\begin{center}
  {\arabicfont \Large
    \textarabic{عَلِّمْ - عَلِّمِي - لَا تُعَلِّمْ - لَا تُعَلِّمِي}
  }
\end{center}

\vspace{0.7cm}

% ─── Practice instruction ───────────────────────────────────
\begin{center}
  নীচের {\arabicfont \textarabic{مصدر}} গুলো থেকে একই নিয়মে
  {\arabicfont \textarabic{فعل}} তৈরী ও মুখস্থ করো।
\end{center}

\vspace{0.5cm}

% ─── Masdar practice table ──────────────────────────────────
\setlength{\arrayrulewidth}{0.5pt}
\setlength{\tabcolsep}{4pt}
\renewcommand{\arraystretch}{2.0}

\noindent\begin{tabularx}{\textwidth}{|C|C|C|C|}
  \hline
  \rowcolor{lightgreen}
  আলোকিত করা & {\arabicfont \textarabic{التَّنْوِيرُ}} &
  পরিষ্কার করা & {\arabicfont \textarabic{التَّنْظِيفُ}} \\
  \hline

  কথা বলা & {\arabicfont \textarabic{التَّكْلِيمُ}} &
  সালাম দেয়া & {\arabicfont \textarabic{التَّسْلِيمُ}} \\
  \hline

  \rowcolor{lightgreen}
  মিথ্যা মনে করা & {\arabicfont \textarabic{التَّكْذِيبُ}} &
  সত্য মনে করা & {\arabicfont \textarabic{التَّصْدِيقُ}} \\
  \hline
\end{tabularx}

